Sign language is the primary communication mode for over 70 million deaf and 
hard-of-hearing people worldwide \cite{who2021}. ASL specifically — used mostly 
in North America — uses hand gestures and finger-spelling to express words and 
concepts. Despite how widely it's used, there are almost no real-time automated 
translation tools available. That gap makes everyday communication with the 
hearing majority significantly harder than it needs to be.

CNNs and real-time object detection frameworks have opened up new options for 
gesture recognition. Systems that can classify a gesture from a live video frame 
in milliseconds are a meaningful step toward accessible communication technology. 
Our approach uses two stages: first, a MediaPipe hand tracking module localizes 
the hand in each camera frame, producing a tight bounding box crop (not visible 
to the user) and a skeleton overlay (visible). Second, a fine-tuned ResNet-50 
\cite{he2016} classifies the cropped image into one of 26 ASL alphabet letters. 
The result is overlaid on the live stream in real time.

Our contributions:
\begin{itemize}
    \item A modular training pipeline with configurable optimizer, scheduler, and early stopping, implemented in PyTorch.
    \item MediaPipe-based hand detection integrated with a ResNet-50 classification head for end-to-end ASL letter recognition.
    \item A real-time inference module with webcam integration and visual feedback.
    \item A FastAPI web interface for browser-based interaction with the live pipeline.
\end{itemize}